\documentclass[11pt]{article}
\usepackage[margin=1in]{geometry}
\usepackage{amsmath,amssymb}
\usepackage[T1]{fontenc}
\usepackage{xcolor}
\usepackage[hidelinks]{hyperref}
\usepackage{needspace}

\setlength{\parindent}{0pt}
\setlength{\parskip}{0.65em}
\setlength{\emergencystretch}{2em}
\newcommand{\status}[1]{\textbf{#1}}

\title{Technical Review of ``Early-Stopping Activation Steering for Factual Preference Alignment in Vietnamese Medical Language Models''}
\author{}
\date{September 4, 2026}

\begin{document}
\maketitle

\section*{Overall assessment}

The manuscript studies a potentially useful inference-time intervention and reports paired results on a fixed set of 500 Vietnamese medical questions.  The schedule arithmetic, category percentages, matched-dose coefficient, and the exact McNemar tests reported from the displayed contingency tables are internally consistent.  The paper also correctly acknowledges that its reference-preference score is not clinician-adjudicated factual accuracy.

The manuscript nevertheless requires major revision before its central conclusions are supported.  The primary endpoint cannot establish clinical factuality, the activation-norm experiment is internally inconsistent with the stated fixed teacher-forcing protocol, important benchmark and implementation details are absent, the label-shuffle control does not constitute a convincing null as described, and the discussion relies on an unreported ``Experiment~10.''  Several strong claims---including prevention of magnitude growth, elimination of repetition, clinical applicability, and superiority to Hybrid RAG---must either be supported by additional evidence or narrowed.

I explicitly checked the steering equations, cumulative-dose derivation, sign conventions, symbol definitions, layer/index terminology, dimensional consistency, branch conventions, tables, cross-section claims, citations, and reference-list hygiene.  The dose derivation in Eqs.~(5)--(7) is correct for a fixed realized horizon $T$ and positive source value $\alpha_0$.  No inverse-trigonometric or coordinate-transform expressions occur, so there is no branch or quadrant convention to repair.  All citation keys used in the text have matching bibliography entries, and no bibliography-wide duplicated-punctuation defect was found.

\section*{Highest-priority fixes}

\begin{enumerate}
    \item Resolve and rerun the activation-norm analysis with an unambiguous shared-context or free-generation protocol; remove the current causal claim about magnitude growth until the reported tensors are consistent with that protocol.
    \item Reframe the outcome as semantic reference preference unless clinician adjudication and category-specific factual checks are added.  Dosage numbers, units, negation, contraindications, and interaction direction require direct evaluation.
    \item Document benchmark construction, counter-answer generation, expert quality control, group-disjoint splitting, and selection of the 500-question subset sufficiently to exclude leakage and selection bias.
    \item Add complete implementation details for activation extraction, batching/padding, hook behavior, BERTScore truncation, RAG construction, and timing; release code and per-question outcomes if possible.
    \item Either report the repetition-penalty experiment in Methods and Results with the primary endpoint, or remove all ``Experiment~10'' and ``resolves completely'' claims.
\end{enumerate}

\section*{Technical and scientific findings}

\subsection*{1. Activation-norm evidence conflicts with the stated protocol}

\status{Classification: definite internal inconsistency.}

\status{Location.} Introduction contribution list (\texttt{paper.tex:55}); hook specification (lines 128--135); ``Empirical Activation-Norm Dynamics'' (lines 220--221).

\status{Problem.} The contribution list says the norms were measured ``under fixed teacher-forcing context.''  The hook is attached to the output of \texttt{model.model.layers[8]} and is the identity for $t>K$.  With a common teacher-forced token sequence, the inputs and attention cache up through block index~8 are not changed by an output hook on that same block.  Consequently, at $t=50>K$, the reported post-hook Layer-8 activation should coincide with the baseline for both Hard Cutoff and Linear Decay.  The manuscript instead reports 53.39 for Linear Decay, 52.75 for Hard Cutoff, and 53.39 for baseline.  If each condition follows its own generated sequence, different norms are possible, but then the context is not fixed and token-content effects are confounded with the intervention.

The interpretation also exceeds the displayed evidence.  Continuous steering changes from 56.28 at $t=10$ to 56.34 at $t=50$, which shows a persistent offset, not progressive ``magnitude inflation'' or growth.  The phrase ``a mild norm undershoot of $0.00$ units'' is self-contradictory.

\status{Why it matters.} Activation relaxation is presented as mechanistic evidence for the proposed method.  Under the stated hook location, the equality after $K$ may be guaranteed by the measurement design rather than discovered empirically, while the Hard Cutoff value suggests that a different protocol was actually used.

\status{Suggested fix.} State whether all conditions use one identical teacher-forced sequence or their own generated tokens.  Specify whether the recorded tensor is the pre-hook or post-hook output, how caches are handled, and how norms are aggregated.  Under shared teacher forcing, verify exact per-prompt equality at block index~8 after $K$ and instead measure downstream blocks, logits, or cached representations if persistence is the target.  Under free generation, call the analysis trajectory-dependent and separate token-content effects.  Report curves with sample sizes and uncertainty, and replace ``growth'' with ``persistent elevation'' unless a positive time trend is demonstrated.

\subsection*{2. Reference preference does not establish factual or clinical correctness}

\status{Classification: unsupported central claim.}

\status{Location.} Abstract (lines 31--34), Introduction (lines 43--49), BERTScore criterion (lines 116--121), Results (lines 164--186 and 223--233), and Conclusion (line 332).

\status{Problem.} The binary endpoint only asks whether BERTScore is larger against $y_i^+$ than against one synthetic $y_i^-$.  This neither verifies the facts in the generated answer nor detects a third error absent from both references.  The benchmark specifically targets dosage, pregnancy contraindication, and interaction errors, where a number, unit, negation, or direction can reverse clinical meaning while leaving embedding similarity high.  The manuscript acknowledges the proxy limitation but repeatedly converts the result into ``factual accuracy,'' ``hallucination mitigation,'' ``clinical safety,'' and evidence for clinical deployment.

\status{Why it matters.} The reported 4.20--5.20 percentage-point gains may reflect movement toward the wording of $y^+$ rather than fewer medical errors.  This is the main validity threat to the paper's title, abstract, and conclusion.

\status{Suggested fix.} Until direct validation is available, consistently use ``BERTScore reference preference'' and avoid ``accuracy'' unless qualified as reference-preference accuracy.  Add blinded clinician adjudication on a prespecified sample and report inter-rater agreement.  Add category-specific checks for medication identity, contraindication polarity, interacting agents and mechanism/direction, dose value, unit, body-weight normalization, frequency, and patient subgroup.  Validate RefPref against those labels and report sensitivity, specificity, and failure examples.

\subsection*{3. Benchmark construction and split integrity are not reproducible}

\status{Classification: definite omission and likely leakage/selection risk.}

\status{Location.} Contribution~1 (line 53) and ``ViHaluEval-Medical Benchmark Construction'' (lines 74--75).

\status{Problem.} The three stated category counts, $165+160+175=500$, describe the evaluation subset, but the sentence grammatically attaches them to the 14,700-record benchmark.  No category counts are given for the full data or the 10,290/2,205/2,205 splits.  ``Expert-structured'' is unsupported because the experts, instructions, adjudication, and quality controls are not described.  The method used to synthesize $y_i^-$ is absent.  ``Question-disjoint'' is stated only for the 500-question evaluation subset; the manuscript does not establish group-disjoint train/validation/test construction or protection against repeated drugs, source passages, templates, paraphrases, and answer pairs.  The rule and seed used to select 500 of 2,205 test pairs are also missing.  Finally, the RAG corpus contains 14,576 passages while the benchmark contains 14,700 pairs, with no mapping or exclusion accounting.

\status{Why it matters.} Template or source overlap can inflate both the learned direction and evaluation score.  Unreported subset selection can also bias the primary results, and the benchmark cannot be independently assessed or reproduced.

\status{Suggested fix.} Add a dataset-flow table showing source units, exclusions, deduplication, category counts, and split counts.  Define the grouping key used before splitting, and report overlap audits at question, drug, source-passage, and template levels.  Describe counter-answer generation and expert review, including qualifications, number of annotators, agreement/adjudication, and acceptance criteria.  State the 500-item sampling rule and seed.  Explain the 14,700-versus-14,576 difference and provide stable dataset/code access plus the source-license analysis.

\subsection*{4. Activation extraction and hook implementation remain ambiguous}

\status{Classification: likely implementation-affecting issue.}

\status{Location.} Eqs.~(1)--(2) and surrounding text (lines 77--86), Eqs.~(3)--(4) (lines 94--106), and implementation protocol (lines 123--139).

\status{Problem.} In Eq.~(1), \texttt{[:, -1, :]} may select padding for shorter examples in a right-padded batch.  The paper does not state the padding side, batching policy, attention-mask-based final index, inclusion of EOS, chat template, prompt/answer separators, truncation, or whether positive and negative responses are length/style matched.  A final-token contrast can therefore encode answer length, terminal punctuation, EOS behavior, or template artifacts rather than truthfulness.  The $N$ in $\mu_l^\pm$ is not explicitly identified as $N_{\mathrm{train}}=10{,}290$ and is later overloaded for prompts and random vectors.

There is also required terminology drift: the abstract and equations call the target ``Layer~8,'' whereas \texttt{model.model.layers[8]} is explicitly the zero-indexed ninth decoder block.  The mathematical label $l=8$ and the verbal label ``Layer~8'' therefore have two plausible numbering conventions.

\status{Why it matters.} These choices can materially change $v_{\text{steer}}$, make replication fail, and cause an implementation to intervene at the wrong block.

\status{Suggested fix.} Provide executable pseudocode.  Select the last non-padding response token via the attention mask, state whether EOS is included, and report padding/truncation/chat-template details.  Add controls for answer length and surface form or pool over aligned answer tokens.  Use $N_{\mathrm{train}}$ in Eq.~(2).  Replace every ambiguous ``Layer~8'' with a single convention such as ``decoder block index~8 (the ninth block),'' including all $h_8^{(t)}$ descriptors.

\subsection*{5. The label-shuffled direction is not a convincing null as described}

\status{Classification: likely invalid control.}

\status{Location.} ``Contrastive Vector Estimation \& Directional Controls'' (lines 88--92) and Table~V (lines 239--261).

\status{Problem.} The manuscript says labels are swapped across $N_{\mathrm{flipped}}=185$ training pairs.  Relative to 10,290 training pairs, this is only about 1.8\%.  If all remaining labels retain their original sign, the resulting vector is expected to preserve most of the original signal and is not a label-shuffle null.  The reported cosine of 0.5839 with $v_{\text{steer}}$ reinforces that the control retains substantial alignment.  The section also says ``two controls,'' then repeats Negative Steering inside a second item that lists four controls.

\status{Why it matters.} The direction-specificity claim relies on excluding generic or label-independent perturbations.  A mostly unshuffled direction does not perform that test.

\status{Suggested fix.} Clarify what 185 counts.  For a null, independently permute all labels or assign balanced random signs to all training differences, repeat this procedure over many seeds, and report the distribution of cosine similarity and test performance.  Keep covariance-matched and isotropic controls separate, and remove the duplicated Negative Steering description.

\subsection*{6. Long-generation BERTScore may be truncated and metric definitions are incomplete}

\status{Classification: likely issue requiring an artifact check.}

\status{Location.} Metric definition (lines 116--121), high-cap evaluation (lines 143--162), ablation table (lines 286--313), and reproducibility paragraph (lines 138--139).

\status{Problem.} BERTScore~0.3.13 tokenizes with the evaluator model's maximum length and truncation enabled; its own documentation warns that BERT-family inputs beyond 510 content tokens are truncated.  The selected \texttt{bert-base-multilingual-cased} tokenizer has a 512-token limit, while the natural-completion experiment permits 800 new model tokens.  The manuscript gives no evaluator-token length distribution or number of truncated outputs.  If truncation occurs, later hallucinations and repetitions do not affect RefPref or mean BERTScore even though Rep-4 is computed on the full output.

The paper also omits BERTScore IDF/rescaling settings and hash, normalization, empty-output handling, and tie precision.  ROUGE-L, Rep-4, EOS Hit, latency, Recall@$k$, and MRR lack complete definitions and aggregation rules.

\status{Why it matters.} Truncation can change the primary high-cap ranking and makes the endpoint cover a different portion of the output than the repetition metric.  Missing settings prevent exact reproduction.

\status{Suggested fix.} Report generated lengths in both Qwen and mBERT wordpieces, the number and condition-wise rate of BERTScore truncation, and a sensitivity analysis using a prespecified common evaluation window or a validated long-text/chunked procedure.  Report the BERTScore hash and all options.  Formally define every metric, including tokenization, treatment of outputs shorter than four tokens, whether EOS is checked before the cap, and whether latency includes prefill, decoding, retrieval, and synchronization.

\subsection*{7. Hyperparameter and schedule ablations omit the primary endpoint}

\status{Classification: unsupported ablation conclusion.}

\status{Location.} Contribution~6 (line 58), factorial ablation (lines 281--313), and layer-wise probing (line 325).

\status{Problem.} Table~VII calls itself a controlled $3\times3$ ablation but does not report RefPref, the paper's primary endpoint, or an unsteered baseline.  The schedule comparisons are therefore interpreted from BERTScore differences as small as 0.0001--0.0018, without paired uncertainty.  The manuscript also does not show the validation procedure that selected block index~8, $K=16$, or $\alpha_0=18$, nor whether all choices were fixed before the test set was evaluated.  The statement that a narrow layer range ``confirms framework robustness'' is stronger than a 100-example validation sweep supports.

\status{Why it matters.} The paper cannot establish that Linear Decay or Hard Cutoff is robustly preferable, or that the reported test comparison was not selected after inspecting multiple configurations.

\status{Suggested fix.} Add RefPref counts and paired differences for every cell, plus baseline rows and uncertainty.  Publish the validation grid, selection metric, tie-breaking rule, seeds, and chronology; reserve the test set for the final locked configuration.  Describe the layer result as limited sensitivity on the sampled validation set unless replicated with uncertainty.

\subsection*{8. The discussion relies on an experiment absent from the manuscript}

\status{Classification: definite unsupported claim.}

\status{Location.} ``Multi-Metric Pareto Trade-off \& Deployment Mitigation'' (line 329).

\status{Problem.} ``Experiment~10'' is not defined in Methods, Results, a table, or an appendix.  Only Rep-4 and latency numbers are quoted.  No RefPref, BERTScore, EOS, output length, contingency table, or statistical uncertainty is supplied for $\theta_{\mathrm{rep}}=1.15$.  The text nevertheless says the penalty ``resolves this trade-off completely'' and later suggests it can ``preserve the +5.20~pp RefPref gain.''  The latter outcome is not reported.  The same natural-completion trade-off is then described a second time in the same paragraph.

\status{Why it matters.} This is the paper's proposed remedy for its clearest adverse result, but readers cannot verify that factual preference is retained or that the latency reduction is not caused by substantially shorter answers.

\status{Suggested fix.} Add the experiment to Methods and Results with the same paired 500 questions and all primary/secondary metrics, including output length and EOS.  Compare it directly with baseline and unpenalized Hard Cutoff.  Otherwise delete the experiment-specific numbers and state repetition-penalty testing only as future work.

\needspace{6\baselineskip}
\subsection*{9. RAG comparisons are under-specified and partly overstated}

\status{Classification: likely comparability problem and unsupported wording.}

\status{Location.} RAG results and Tables~V--VI (lines 223--279), Discussion (line 329), and Conclusion (line 332).

\status{Problem.} BM25, BGE-M3, Hybrid RRF, and oracle RAG appear only in Results.  The manuscript does not specify passage construction, dense-model revision and pooling, index normalization, retrieved $k$, RRF candidate depth and exact formula, prompt template, context truncation, or whether retrieval time is included.  The claim that steering ``outperforms Hybrid RAG'' is not supported inferentially by the displayed comparison ($p=0.5682$); the defensible statement is that its point estimate is 1.00~pp higher and statistically indistinguishable on this sample.  ``Zero context memory'' is also too absolute: steering avoids retrieved-context KV-cache cost but still stores a vector and runs a hook.

All non-RAG methods are assigned the identical $7.32\pm0.84$~s in Table~V.  Without condition-specific timing and a paired timing analysis, this looks like a reused baseline estimate rather than evidence that the hook adds no measurable cost.

\status{Why it matters.} Retrieval quality and latency depend strongly on these choices.  The current description does not support a fair systems comparison or the zero-overhead language.

\status{Suggested fix.} Add a RAG methods subsection and release the corpus/index configuration.  State timing boundaries, batch size, warm-up, synchronization, hardware allocation, and whether retrieval/index transfer is included.  Report condition-specific paired latency differences and confidence intervals.  Use ``numerically higher but not significantly different'' for Hybrid RAG, and ``no retrieved-context overhead'' instead of ``zero context memory.''

\subsection*{10. Matched-dose and statistical interpretations need tighter case definitions}

\status{Classification: likely issue; displayed arithmetic otherwise checks out.}

\status{Location.} Eqs.~(5)--(7) and matched-dose paragraph (lines 108--114), Tables~III--VII, and random-control interpretation (line 226).

\status{Problem.} The matched-dose equality assumes a fixed realized $T=200$.  Table~VII reports a nonzero EOS rate for one cell, and the general method permits early termination; for a sequence ending before $T$, a constant coefficient based on 200 steps no longer has dose $D_{\mathrm{decay}}$.  The manuscript needs a case distinction between maximum horizon and realized length.  Also, the schedule called ``linear decay'' remains at $\alpha_0/K$ at $t=K$ and then jumps to zero; the equation is correct as written, but this is linear decay followed by a cutoff, not a linear ramp reaching zero at $K$.

The displayed exact McNemar $p$-values agree with the contingency counts.  However, the Monte Carlo value $1/101$ is only a comparison with 100 selected random directions and has that value as its minimum attainable plus-one estimate.  Treating the learned direction as exchangeable with isotropic directions requires a stated null construction; the current result is descriptive rather than confirmation of factual specificity.  Multiple comparisons extend beyond the three natural-completion tests covered by Holm adjustment.

\status{Why it matters.} Ambiguous horizons undermine the claimed dose control, while the current random-direction and multiplicity language overstates inferential strength.

\status{Suggested fix.} Define $T$ as maximum or realized length and write the dose using $\min(T,K)$ where necessary; either force every matched-dose run to 200 steps or match dose per realized sequence.  Rename the schedule or use denominator $K-1$ if zero at $t=K$ is intended.  Prespecify primary comparisons, define the Monte Carlo null, increase the number of random/shuffled draws, and report adjusted or explicitly exploratory secondary analyses.

\section*{Meaningful editorial and presentation findings}

\begin{enumerate}
    \item \status{Location: lines 138--139.} The hardware/software environment sentence is duplicated almost verbatim.  Merge it into one paragraph and retain the random-seed details once.
    \item \status{Location: lines 88--92.} ``We evaluate two controls'' conflicts with a list that repeats Negative Steering and then names four conditions.  State the exact number of unique controls and give each one item.
    \item \status{Location: lines 220--221.} Replace ``a mild norm undershoot of $0.00$ units'' with ``no difference at the reported precision'' if the value is genuinely equal; use ``undershoot of $0.64$'' only for the 52.75 versus 53.39 comparison and only after resolving the protocol.
    \item \status{Location: lines 223--233.} BM25 performance is described twice in adjacent paragraphs.  Consolidate the retrieval results, then present the paired comparisons once.
    \item \status{Location: line 329.} The natural-completion trade-off is repeated within one very long paragraph.  Retain one concise statement, and separate observed results from proposed deployment mitigations.
    \item \status{Location: throughout.} Use one capitalization and term for ``early-stopping steering''; avoid the late variant ``Early-Stop Steering.''  Likewise, capitalize generic ``small language models'' only when beginning a sentence or defining the abbreviation.
\end{enumerate}

\section*{Proofing-sweep items}

The supplied mechanical scan was run on \texttt{paper.pdf}.  Its two hits were false positives: the apparent double period is the mathematical ellipsis in $\{1,\ldots,100\}$, and the semicolon-capitalization hit precedes the metric name ``Recall@3.''  Manual spot-checking identified these high-confidence corrections:

\begin{itemize}
    \item \status{Lines 138--139:} delete the duplicated reproducibility sentence.
    \item \status{Line 221:} correct the contradictory phrase ``norm undershoot of $0.00$ units.''
    \item \status{Line 233:} typeset ``+7.13s'' as ``$+7.13$~s'' for consistent number--unit spacing.
    \item \status{Line 329:} remove the duplicated trade-off passage and the broken internal reference to ``Experiment~10'' unless that experiment is added and labeled.
    \item \status{Table~VII:} do not bold both 0.7151 and 0.7150 as the unique maximum unless the intent is to mark values tied at a declared precision; state that rule in the caption.
\end{itemize}

\section*{Items requiring artifact or external verification}

\begin{enumerate}
    \item The reported benchmark, generated answers, per-question RefPref outcomes, bootstrap samples, activation traces, and GPU logs are not included in the manuscript.  The numerical tables cannot be independently reproduced from the available files.
    \item The stated Tesla T4 plus \texttt{bfloat16} compute configuration needs a runtime record.  NVIDIA's published T4 precision specifications enumerate FP32, FP16, INT8, and INT4 rather than native BF16.  Report \texttt{torch.cuda.is\_bf16\_supported()}, the actual \texttt{Linear4bit.compute\_dtype}, any fallback/casting behavior, and the dtype of the hooked tensor.
    \item Verify that every citation supports the exact nearby claim.  In particular, the specific medical-error examples attributed near \texttt{ref3} and the ``internal activation pathways'' explanation attached to \texttt{ref11} are more specific than the manuscript's summaries of those sources establish.
    \item Verify the official English/transliterated title and publisher formatting for \texttt{ref6}, and cite BGE-M3 and the exact dense-retrieval model revision used in the RAG experiments.
\end{enumerate}

\section*{Recommended decision}

Major revision.  The paired preference results are potentially publishable, and their displayed arithmetic is largely sound, but the current manuscript does not yet establish reduction of medical hallucinations or the proposed activation-norm mechanism.  A revised paper should make semantic reference preference the bounded claim, repair the norm experiment, document the benchmark and implementation, and either supply or remove the unreported repetition-mitigation result.

\end{document}